\documentclass[12pt,a4paper,fontset=none]{ctexart}

% ---------- 页面 ----------
\usepackage{geometry}
\geometry{top=2.8cm,bottom=2.8cm,left=2.8cm,right=2.8cm}

% ---------- 字体（宋体正文 / 黑体标题 / 楷体引文）----------
\setCJKmainfont{AR PL UMing CN}
\setCJKsansfont{WenQuanYi Zen Hei}
\newCJKfontfamily\kai{AR PL KaitiM GB}
\newCJKfontfamily\hei{WenQuanYi Zen Hei}

% ---------- 颜色 ----------
\usepackage{xcolor}
\definecolor{ink}{HTML}{1F2D3D}
\definecolor{accent}{HTML}{8C2A2A}
\definecolor{accentbg}{HTML}{FBF1EE}
\definecolor{notebg}{HTML}{F3F5F7}
\definecolor{notebar}{HTML}{9AA7B4}
\definecolor{link}{HTML}{2C5F7C}
\definecolor{rule}{HTML}{C9A24B}

% ---------- 排版增强 ----------
\usepackage[most]{tcolorbox}
\usepackage{enumitem}
\usepackage{microtype}
\usepackage{titlesec}
\usepackage{fancyhdr}
\usepackage[hidelinks]{hyperref}

\linespread{1.5}
\setlength{\parskip}{0.45em}
\setlength{\parindent}{2em}

% ---------- 章节样式 ----------
\ctexset{
  section/format = \Large\hei\bfseries\color{accent},
  section/beforeskip = 2.2ex,
  section/afterskip = 1.2ex,
  subsection/format = \large\hei\bfseries\color{ink},
  subsection/beforeskip = 1.6ex,
  subsection/afterskip = 0.8ex,
}

% ---------- 页眉页脚 ----------
\pagestyle{fancy}
\fancyhf{}
\renewcommand{\headrulewidth}{0.2pt}
\fancyhead[L]{\small\hei\color{notebar}哲学的故事}
\fancyhead[R]{\small\hei\color{notebar}一场跨越2500年的对话}
\fancyfoot[C]{\small\color{notebar}\thepage}

% ---------- 自定义环境 ----------
% 主线悬念框
\newtcolorbox{mainline}{
  enhanced, breakable,
  colback=accentbg, colframe=accent, boxrule=0.6pt,
  left=12pt,right=12pt,top=9pt,bottom=9pt,
  fontupper=\kai\color{ink}, arc=2pt,
  borderline west={3pt}{0pt}{accent}
}
% 旁注 / 提示框
\newtcolorbox{note}{
  enhanced, breakable,
  colback=notebg, colframe=notebg, boxrule=0pt,
  left=12pt,right=12pt,top=7pt,bottom=7pt,
  fontupper=\kai\color{ink!85}, arc=1pt,
  borderline west={3pt}{0pt}{notebar}
}
% 延伸阅读
\newcommand{\link}[1]{{\par\smallskip\noindent\hspace*{2em}\small\sffamily\color{link}\,»\,#1\par\smallskip}}
% 名言强调
\newcommand{\saying}[1]{{\kai\color{accent}#1}}

\begin{document}

% ================= 封面 =================
\thispagestyle{empty}
\begin{center}
\vspace*{3.2cm}
{\color{rule}\rule{0.5\textwidth}{0.8pt}}\\[1.1cm]
{\Huge\hei\bfseries\color{ink} 哲 学 的 故 事}\\[0.7cm]
{\LARGE\kai\color{accent} 一场跨越 2500 年的对话}\\[1.1cm]
{\color{rule}\rule{0.5\textwidth}{0.8pt}}\\[2.2cm]
{\large\kai\color{ink!80} 人能认识世界吗？人该怎样活？\\[0.4em]
而到了最后——我们造出的机器，能思考吗？}\\[3.5cm]
{\normalsize\sffamily\color{notebar} 哲学体系 · 叙事式入口}\\[0.3em]
{\small\sffamily\color{notebar} 2026}
\end{center}
\clearpage

% ================= 引言 =================
\thispagestyle{fancy}
\begin{note}
这不是教科书，是一个\textbf{故事}。把全库的内容（西方、东方、所有专题）当成\textbf{同一场对话}来讲：每个人都在回答前一个人留下的问题，又给后人留下新的难题。读完它，你会拿到整座知识库的\textbf{地图与情节线}。
\end{note}

\begin{mainline}
主线只有一个悬念：\textbf{人能认识世界吗？人该怎样活？而到了最后——我们造出的机器，能思考吗？}
\end{mainline}

% ================= 序幕 =================
\section*{序幕\quad 第一个抬头问“为什么”的人}
\addcontentsline{toc}{section}{序幕}

很久很久以前，人类用神话解释一切：打雷是神发怒，季节是诸神的故事。直到大约公元前 6 世纪，在地中海、黄河、恒河三处地方，几乎\textbf{同时}，有人不再满足于“神说”，开始问：\textbf{事物本身，到底是怎么回事？}

哲学家雅斯贝尔斯把这段时间叫“轴心时代”。故事就从这里，分三条河同时开始流。

% ================= 第一卷 =================
\section{第一卷 · 西方：理性的冒险}

\subsection{第一章\quad 米利都的海边——世界由什么构成}

故事的第一个主角叫\textbf{泰勒斯}。他站在米利都的海边，没有求神，而是大胆断言：\textbf{万物的本原是水。}答案错了，但提问的方式对了——这是人类第一次试图用\textbf{一个自然原因}解释整个世界。哲学诞生了。

接力棒传下去：阿那克西曼德说本原是“无限定者”，赫拉克利特说\saying{“人不能两次踏进同一条河”}——世界是永恒的流变与斗争；巴门尼德针锋相对：\saying{“存在者存在，不存在者不存在”}——变化是幻觉，真实是不变的“一”。

\begin{note}
第一个大裂缝出现了：世界的本质是\textbf{流变}（赫拉克利特）还是\textbf{不变}（巴门尼德）？这道裂缝，两千年后还在。
\end{note}
\link{细节见《古希腊哲学》}

\subsection{第二章\quad 苏格拉底转身——把目光从星空转向人心}

正当大家争论自然，雅典街头出现一个其貌不扬、爱缠着人问问题的老头\textbf{苏格拉底}。他做了一件惊天动地的事：\textbf{把哲学从“世界是什么”转向“人该怎么活”。}

他从不给答案，只反问——“你说正义是什么？”“你确定吗？”——直到对方发现自己其实什么都不懂。\saying{“我唯一知道的，就是我一无所知。”}最后雅典人受不了他，判他死刑。他平静地喝下毒酒，用死亡证明：有些东西比生命更重要。

他的学生\textbf{柏拉图}痛心之余，写下一个流传千古的画面——\textbf{洞穴}：我们都是被锁在洞里的囚徒，把墙上的影子当成真实。哲学家就是那个挣脱锁链、走出洞口、看见真正太阳的人。真实的世界不在感官里，而在永恒的“理念”中。

柏拉图的学生\textbf{亚里士多德}却把目光拉回地面：“吾爱吾师，吾更爱真理。”他不信天上的理念，要研究眼前每一个具体事物——他几乎一个人开创了逻辑学、物理学、生物学、伦理学。他说幸福不是快感，而是\textbf{德性的实现}，是“在恰当的时间、对恰当的人、做恰当的事”。

\begin{note}
师徒三代立起了西方哲学的两座地基：柏拉图的\textbf{理念}（真理在彼岸）与亚里士多德的\textbf{经验}（真理在此岸）。后世几乎所有争论，都是这两条路的延续。
\end{note}

\subsection{第三章\quad 帝国黄昏——如何在乱世中安顿自己}

亚历山大的帝国崩塌，城邦瓦解，人心惶惶。哲学这时变成了\textbf{心灵的药方}：

\begin{itemize}[leftmargin=2.4em,itemsep=0.2em]
\item \textbf{斯多葛派}说：分清“能控制的”和“不能控制的”，只为前者烦恼。命运无法改变，但我面对它的态度可以。
\item \textbf{伊壁鸠鲁派}说：快乐是善，但最高的快乐是\textbf{宁静}——免于恐惧和欲望的折磨。
\item \textbf{怀疑派}说：那就别下判断了，悬置一切，自然就平静了。
\end{itemize}

\begin{note}
这是哲学第一次大规模“落地”为生活方式。两千年后它会以“斯多葛晨间练习”的名字回到你的日程表。
\end{note}

\subsection{第四章\quad 上帝的千年——信仰与理性能否两全}

罗马拥抱了基督教，哲学的主题变成：\textbf{信仰与理性，如何共处？}

\textbf{奥古斯丁}把柏拉图请进教堂——真理之光来自上帝；他在《忏悔录》里剖开自己的灵魂，问“时间究竟是什么”。一千年后，\textbf{阿奎那}把刚从阿拉伯世界传回的亚里士多德请进来，用五种方式“证明”上帝存在，搭起信仰与理性的宏伟桥梁。

别忘了这座桥的另一半工匠是\textbf{伊斯兰与犹太哲学家}：伊本·西那（阿维森纳）、伊本·路世德（阿威罗伊）保存并发展了亚里士多德，迈蒙尼德调和了信仰与哲学。是他们把希腊的火种又传回了欧洲。

末尾，\textbf{奥卡姆}举起一把剃刀：\saying{“如无必要，勿增实体”}——别用复杂的解释，简单的就够了。这把剃刀，今天还躺在每个科学家和工程师的工具箱里。

\subsection{第五章\quad 我思故我在——近代的怀疑与觉醒}

文艺复兴的曙光里，一个法国人\textbf{笛卡尔}决定把一切推倒重来。他怀疑一切：感官会骗人，也许整个世界都是恶魔编织的梦。可是——\textbf{当我在怀疑时，“我正在怀疑”这件事不可怀疑}。\saying{“我思故我在。”}他在怀疑的废墟上，找到了第一块不可动摇的基石。

但他留下一个烫手的难题：既然心灵和身体是两种完全不同的东西，它们\textbf{怎么互动}？（记住这个“心身问题”——故事的最后它会变成“机器能有心灵吗”。）

接下来，欧洲分成两大阵营，吵了一百年：

\begin{itemize}[leftmargin=2.4em,itemsep=0.2em]
\item \textbf{理性主义}（笛卡尔、斯宾诺莎、莱布尼茨）：真知识来自\textbf{理性}。斯宾诺莎甚至说，上帝就是自然，万物都是这唯一实体的两面。
\item \textbf{经验主义}（洛克、贝克莱、休谟）：真知识来自\textbf{经验}。洛克说人心生来是白板。而\textbf{休谟}抛出最致命的一问——\saying{“你凭什么相信明天太阳还会升起？”}你见过它升起一万次，但“一万次”逻辑上推不出“第一万零一次”。\textbf{因果，也许只是习惯。}
\end{itemize}

\begin{note}
休谟这一问，叫“归纳问题”。它像一颗定时炸弹，两百年后会在科学哲学和机器学习里再次引爆。
\end{note}

同时，另一群人在问\textbf{人该如何共同生活}：霍布斯说自然状态是“一切人对一切人的战争”，所以需要“利维坦”；洛克说人有天赋权利，政府是为保护它而存在；卢梭说“人生而自由，却无往不在枷锁之中”。\textbf{社会契约论}就此奠定了现代政治的地基。

\subsection{第六章\quad 康德的哥白尼革命——是世界绕着我们转}

休谟的炸弹惊醒了一个德国人\textbf{康德}——“是他打断了我的独断论迷梦。”康德憋出了人类思想史上最艰深也最伟大的一招\textbf{哥白尼式革命}：

\begin{mainline}
不是我们的认识符合对象，而是\textbf{对象符合我们的认识}。时间、空间、因果，不是世界本来的样子，而是我们大脑加在世界上的“眼镜”。我们永远看不到“物自体”，只能看到经过眼镜加工的“现象”。
\end{mainline}

他用三大批判划定了人类理性的边界：能知道什么（《纯粹理性批判》）、该做什么（《实践理性批判》：\saying{“要只按照你愿意它成为普遍法则的准则去行动”}）、可以希望什么（《判断力批判》）。

康德之后，\textbf{黑格尔}把整个历史看成“绝对精神”通过正—反—合不断展开的宏大戏剧；\textbf{马克思}把黑格尔倒过来——“哲学家们只是解释世界，问题在于改变世界”，推动历史的不是精神，是\textbf{物质与经济}；\textbf{叔本华}说世界的本质是盲目的“意志”，人生在痛苦与无聊间摇摆；\textbf{尼采}则喊出那句惊雷——\saying{“上帝死了”}——旧的价值崩塌了，人必须成为自己的立法者，成为“超人”，活出“权力意志”。

\subsection{第七章\quad 二十世纪的两条岔路——语言与存在}

进入现代，哲学分成两条几乎不再对话的大路：

\textbf{分析哲学}这条路，相信问题出在\textbf{语言}上。弗雷格、罗素发明了现代逻辑；天才\textbf{维特根斯坦}先在《逻辑哲学论》里宣布\saying{“凡不能说的，就应当沉默”}，企图一举划清语言的边界；晚年又亲手推翻自己，提出\textbf{“意义即使用”}——词语没有固定本质，意义就活在它被使用的“语言游戏”里。而\textbf{哥德尔}则用一条定理震惊了所有人：\textbf{任何足够强的形式系统，都有它自己证明不了的真理}——理性有它够不到的地方。

\textbf{欧陆哲学}这条路，关心\textbf{存在}本身。胡塞尔要“回到事物本身”（现象学）；他的学生\textbf{海德格尔}追问那个被遗忘了两千年的问题——\textbf{“存在”究竟是什么}，并警告技术正在把万物（包括人）变成可调用的“资源”；\textbf{萨特}喊出\saying{“存在先于本质”}——人没有预先给定的本质，你的选择造就你是谁，因此“人是被判定为自由的”，这自由沉重得令人眩晕；\textbf{加缪}则面对荒诞——世界没有意义，但我们依然要像西西弗斯推石上山那样，\textbf{在反抗中活出意义}。

最后，\textbf{福柯}揭示知识背后是权力，\textbf{德里达}拆解一切文本的确定意义——后现代登场，宏大叙事纷纷瓦解。

西方这条河，从“世界由什么构成”出发，绕了 2500 年，最后流到“语言、存在与权力”。现在，让我们回到起点，看另外两条同样古老的河。

% ================= 第二卷 =================
\section{第二卷 · 东方：另外两条大河}

\subsection{第八章\quad 黄河之畔——儒与道，入世与出世}

就在苏格拉底之前一百年，黄河边上，\textbf{孔子}带着弟子周游列国，到处碰壁，却始终相信：天下无道，更要有人去担当。他不谈玄远的本体，只问\textbf{人与人之间该如何相处}——“仁”是爱人，“礼”是恰当的分寸，\saying{“己所不欲，勿施于人”}。

\textbf{孟子}接棒，坚信人性本善，每个人心里都有“四端”，像火种，扩充开来就是圣贤；\textbf{荀子}偏要唱反调——人性本恶，正因为恶，才更需要后天的礼法去“化性起伪”。这场\textbf{性善 vs 性恶}的辩论，是中国版的人性大问。后来宋明的\textbf{朱熹}讲“存天理”，\textbf{王阳明}讲\saying{“知行合一”}“心即理”，把儒学推向哲学的高峰。

与孔子同时，传说中骑青牛西去的\textbf{老子}留下五千言《道德经》，说的却是完全相反的智慧：\saying{“道可道，非常道。”}最高的道说不出口；最好的治理是“无为”；\saying{“上善若水”}——最强大的恰恰是最柔弱的。

到了\textbf{庄子}，这智慧变成了汪洋恣肆的故事：鲲鹏展翅九万里的“逍遥游”；庖丁解牛“游刃有余”——最高的技艺已超越技艺本身；梦见自己变成蝴蝶，醒来困惑“是我梦见蝴蝶，还是蝴蝶梦见我”。庄子说：\textbf{真正的领会，超越语言}——\saying{“得意而忘言”}。

\begin{note}
记住“得意忘言”——故事结尾，它会变成质问 AI 的一把利刃。儒家教你\textbf{入世担当}，道家教你\textbf{出世逍遥}——这一进一退，成了中国人精神世界呼吸的两半。
\end{note}

\subsection{第九章\quad 恒河之滨——苦的解脱与思辨的森林}

恒河边，一位王子\textbf{乔达摩·悉达多}放弃王位，在菩提树下证悟成佛。他不谈造物主，只诊断人生：\textbf{苦谛}（人生是苦）、\textbf{集谛}（苦因于贪嗔痴）、\textbf{灭谛}（苦可以止息）、\textbf{道谛}（有一条路通向解脱）。核心洞见是“缘起”——万物相互依存，没有任何东西能独立自存，所以\textbf{“无我”}：那个你执着的“我”，其实只是不断变化的因缘聚合。

佛教传入中国，长出新的枝叶：\textbf{禅宗}主张\saying{“不立文字，直指人心，见性成佛”}，六祖慧能那句“本来无一物，何处惹尘埃”成了千古公案。而在印度本土，佛教哲学发展出极其精密的\textbf{唯识学}——“万法唯识”，世界是“识”的变现，连那个执着的“我”，也不过是第七识（末那识）错认的产物。

而印度哲学远不止佛教。那是一片思辨的森林：\textbf{六派正统哲学}（数论、瑜伽、正理、胜论、弥曼差、吠檀多）各成体系，尤其\textbf{因明学}——印度的逻辑学，把“现量”（感知）与“比量”（推理）分得清清楚楚，毫不逊色于亚里士多德。

\begin{note}
三条大河在此并流：希腊问\textbf{世界}，中国问\textbf{社会与人生}，印度问\textbf{苦与解脱}。同一个人类，在同一段时间，问出了完全不同却同样深刻的问题。这本身就是哲学最动人的地方。
\end{note}

% ================= 第三卷 =================
\section{第三卷 · 专题：当大问题落到具体战场}

故事讲到现代，那些古老的大问题分头钻进了各个具体领域，变成今天仍在激烈交锋的专题。

\textbf{科学可靠吗？}\quad \textbf{波普尔}说，科学的标志不是“被证实”，而是\textbf{“可被证伪”}：一个理论必须敢于预言“什么情况下我就错了”。\textbf{库恩}反驳：科学不是平稳积累，而是\textbf{“范式”}的革命性更替，新旧范式之间甚至“不可通约”。拉卡托斯试图调和，费耶阿本德干脆喊“怎么都行”。休谟那颗“归纳”炸弹，在这里第二次引爆。

\textbf{什么是善？}\quad 三大伦理传统各执一词：\textbf{功利主义}说“看结果，最大多数人的最大幸福”；\textbf{义务论}（康德）说“看动机和原则，有些事就是不能做”；\textbf{德性伦理}（亚里士多德）说“别只问该做什么，先问该成为什么样的人”。电车难题、AI 伦理，都是它们的当代战场。

\textbf{如何组织社会？}\quad 从社会契约，到罗尔斯的\textbf{“无知之幕”}（想象你不知道自己将生为谁，你会设计一个怎样的社会？），到社群主义与女权主义的反诘。

\textbf{什么是美？}\quad 康德说审美是“无功利的愉悦”，丹托问“为什么一个普通的盒子放进美术馆就成了艺术”；中国美学则讲“意境”“气韵生动”。

% ================= 终幕 =================
\section*{终幕\quad 所有的河，流向同一个问题}
\addcontentsline{toc}{section}{终幕}

现在，故事来到我们这个时代。

笛卡尔的“心身问题”、图灵的“机器能思考吗”、塞尔的“中文房间”、维特根斯坦的“意义即使用”、庄子的“得意忘言”、唯识的“无我”、休谟的“归纳问题”、亚里士多德的“知道如何 vs 知道那个”——这些跨越 2500 年、来自东西方的所有线索，\textbf{在同一个问题上交汇了}：

\begin{mainline}
我们亲手造出的机器，正在“思考”吗？它“理解”我们的语言吗？如果不算，到底缺了什么？
\end{mainline}

这不再是空想。当 GPT 流畅地与你对话，塞尔会说它只是个巨大的“中文房间”——有语法，没语义；维特根斯坦的信徒会说它学会了“用法”，所以也学会了“意义”；庄子会说它也许是个“数字庖丁”——有出神入化的“how”，却没有可以言说的“that”；唯识学会反问——你凭什么觉得\textbf{你}的理解背后，就一定站着一个真实的“我”？

整座知识库 2500 年的思想，最终都成了你思考这个问题的\textbf{弹药}。而怎样把这些弹药真正打出去——重构论证、构造最强反驳、敢于站队、落到你的 AI 研究上——那是另一种功夫。

% ================= 尾声 =================
\section*{尾声\quad 故事其实没有结尾}
\addcontentsline{toc}{section}{尾声}

哲学最迷人的地方，是 2500 年过去，\textbf{最初的那些问题一个都没被“解决”}——它们只是被问得越来越深、越来越精。

泰勒斯问“世界由什么构成”，今天的物理学家还在问；苏格拉底问“人该怎么活”，今天的你也在问；笛卡尔问“心灵是什么”，今天我们对着自己造的 AI 接着问。

所以读哲学不是为了背下答案，而是为了\textbf{学会被这些问题卷入}——成为这场跨越千年、仍在进行的大对话里，\textbf{下一个开口的人}。

\begin{center}
\vspace{0.6em}
{\kai\color{accent}\Large 故事讲完了。现在，轮到你了。}
\vspace{0.4em}
\end{center}

\vfill
\begin{center}
{\color{rule}\rule{0.35\textwidth}{0.6pt}}\\[0.4em]
{\small\sffamily\color{notebar} 哲学体系 · 叙事式入口 · 把全库串成一个故事}
\end{center}

\end{document}
